%% 編譯器：XeLaTeX（Overleaf: Menu -> Compiler -> XeLaTeX）
\documentclass[12pt]{article}
\usepackage[UTF8]{ctex}
\usepackage[margin=1in]{geometry}
\usepackage{booktabs,array,longtable,enumitem,hyperref}
\onehalfspacing
\newcommand{\path}[1]{\texttt{\detokenize{#1}}}

\title{Manuscript 骨架}
\author{}
\date{}

\begin{document}
\maketitle

給 Overleaf 用的章節骨架。對應 R1--R4 草稿資料夾 $\to$ manuscript 章節；
\textbf{[R2]}、\textbf{[R3]} 這類標記指向已有真實內容、可直接搬過去的地方，不用重寫。

\section*{Title}
\textit{Does the Relation, or Its Change Over Time, Predict Stock Returns?
A Replication of Temporal Relational Ranking with a Dynamic-Relation Extension}（暫定，定稿前再調整）

\section*{Abstract}
一段（150--250 字）：問題 $\to$ 方法 $\to$ 主要發現。等 R3/R4 數字齊了再寫，是全文最後寫的部分。

\section{Introduction}
來源：\path{報告/R1_主題與動機/PLAN.md}，已有草稿，只差潤成正式段落。
\begin{itemize}[leftmargin=*]
  \item 1.1 背景問題：為什麼股票不該被當獨立個體預測
  \item 1.2 本文複製的論文：Feng et al.\ (2019) RSR --- 一段話講完核心方法
  \item 1.3 Research Questions：RQ1（複製）、RQ2（延伸：動態關聯）
  \item 1.4 Contribution 一句話總結
\end{itemize}

\section{Related Work}
新的一節，R1 PLAN.md 沒直接涵蓋，但素材已在手上：
\begin{itemize}[leftmargin=*]
  \item RSR 之外的關係型股價預測方法：HGTAN（hypergraph）、GCNET（graph conv）--- 各一段，
        說明與本文延伸方向（動態 vs.\ 靜態關係圖）的差異
  \item 從 \path{文獻_論文參考/related-literature-shortlist.md} 挑 2--4 篇（供應鏈關聯、跨市場資訊傳導），
        作為「為什麼關聯會隨時間變化」這個假設的外部支持
\end{itemize}

\section{Data}
\textbf{[R2]} 來源：\path{報告/R2_資料與EDA/eda.ipynb}，A1/A2/A3/B1 的發現直接改寫成敘述文字：
\begin{itemize}[leftmargin=*]
  \item 3.1 資料來源：NASDAQ 1,026 檔、\path{data/2013-01-01/}，1,246 個交易日（約 5 年，不是 README 所說的 30 年）
  \item 3.2 關係圖：industry relation，156 檔（15\%）為 ETF/基金、非個股，已排除
  \item 3.3 資料清理：\path{-1234} 缺值代碼、排除 ETF 後的樣本數
\end{itemize}

\section{Methods}
\textbf{[R3]} 三個模型設定，來源：\path{報告/R3_模型與實驗設計/PROGRESS.md}：
\begin{itemize}[leftmargin=*]
  \item 4.1 Baseline：Rank\_LSTM（無關係圖）
  \item 4.2 RSR：固定 industry 關係圖 + Temporal Graph Convolution（複製目標）
  \item 4.3 本文延伸：動態關聯 --- 滾動窗口相關係數，逐窗口重估關係圖
        （\path{preprocess/dynamic_relation.py}，已驗證機制：邊密度隨窗口從 4.6\% 變動到 33.8\%）
\end{itemize}

\section{Experiments and Results}
\begin{itemize}[leftmargin=*]
  \item 5.1 評估指標：MSE、MRR-Top1、模擬報酬（\path{training/evaluator.py}）
  \item 5.2 Baseline 結果 \textbf{[R3 已有真數字]}：Test MSE 0.000377、mrrt 0.049、btl 1.02（NASDAQ 全量，50 epochs）
  \item 5.3 RSR 結果：\textbf{待補}（embedding 已匯出，
        \path{data/pretrain/NASDAQ_rank_lstm_seq-4_unit-32_0.npy}，差最後一次全量訓練）
  \item 5.4 動態關聯結果：\textbf{待補}（腳本可動，尚未接進 \path{relation_rank_lstm.py}、尚未全量跑）
  \item 5.5 Ablation 比較表：baseline / RSR / 動態，三欄同一組指標
\end{itemize}

\section{Discussion}
\textbf{[R4]} 來源：\path{報告/R4_實證分析與結論/PLAN.md} 的三種情境（動態明顯更好 / 沒差異 / 更差），
等 5.4 數字出來後選一種改寫：
\begin{itemize}[leftmargin=*]
  \item 6.1 複製是否成功：與原論文方向是否一致
  \item 6.2 動態關聯的解讀：不管結果好壞都要講原因
  \item 6.3 與原論文的差異：資料期間、市場範圍
  \item 6.4 限制
\end{itemize}

\section{Conclusion and Future Work}
一段總結，並呼應 syllabus 「Potentials of your projects」：若動態關聯有效，下一步可往 thesis/journal publication 延伸。

\section*{References}
Chicago style。\path{文獻_論文參考/} 底下的論文全部要列進來，主要複製論文以 Chicago 格式標明。

\section*{現在能寫的 vs.\ 還要等的}
\begin{center}
\begin{tabular}{ll}
\toprule
章節 & 現在能不能寫 \\
\midrule
1 Introduction & 能寫，素材齊了 \\
2 Related Work & 能寫，素材齊了 \\
3 Data & 能寫，\path{eda.ipynb} 內容直接改寫 \\
4 Methods & 能寫，三個模型設定都定案 \\
5.2 Baseline 結果 & 能寫，真數字已有 \\
5.3 RSR 結果 & 等全量訓練跑完 \\
5.4 動態關聯結果 & 等 \path{relation_rank_lstm.py} 接上動態關係圖並全量跑完 \\
6 Discussion / 7 Conclusion & 等 5.3、5.4 都有數字才能寫 \\
Abstract & 全部寫完最後才寫 \\
\bottomrule
\end{tabular}
\end{center}

\end{document}
